\chapter{Projects\\เล่าโครงการให้เห็นวิธีคิด}

รายชื่อโครงการบอกปริมาณประสบการณ์ แต่เรื่องเล่าโครงการที่ดีเผยให้เห็นวิธีที่ Mentee นิยามปัญหา ตัดสินใจภายใต้ข้อจำกัด ร่วมมือกับผู้อื่น และเรียนรู้จากผลลัพธ์ Projects จึงเป็นพื้นที่ที่องค์ประกอบอื่นของ Tech Portfolio มาบรรจบกัน

\learningobjectives{
  \item เลือกโครงการที่เป็นตัวแทนความสามารถและทิศทาง
  \item เล่าโครงการด้วยตรรกะปัญหา การตัดสินใจ และผลลัพธ์
  \item แยกบทบาทของบุคคลออกจากผลงานของทีมอย่างซื่อตรง
}

\section{เลือก ไม่ใช่รวบรวม}
Portfolio ไม่จำเป็นต้องแสดงทุกโครงการ ควรเลือก 2--4 โครงการที่ช่วยตอบคำถามเชิงกลยุทธ์ เช่น โครงการใดพิสูจน์ความลึก โครงการใดแสดงความสามารถข้ามศาสตร์ โครงการใดเผยการเรียนรู้จากความล้มเหลว และโครงการใดชี้ทิศทางอนาคต

\section{โครงเรื่องห้าจังหวะ}
เรื่องเล่าโครงการที่น่าเชื่อถือประกอบด้วย: \textbf{บริบทและผู้เกี่ยวข้อง} \textbf{ปัญหาที่ค้นพบ} \textbf{การตัดสินใจสำคัญ} \textbf{ผลลัพธ์พร้อมหลักฐาน} และ \textbf{สิ่งที่เรียนรู้} โครงนี้เน้นเหตุผลเบื้องหลังการทำงาน ไม่ใช่เพียงลำดับกิจกรรม

\begin{examplebox}
\textbf{โครงการลด food waste ในตลาดสด}\\
ทีมเริ่มจากสมมติฐานว่าร้านค้าขาดความรู้ แต่การสัมภาษณ์พบว่าคอขวดคือแรงจูงใจทางเศรษฐกิจ Mentee จึงเปลี่ยนจากการอบรมเป็นการทดลองระบบคัดแยกและช่องทางขายใหม่ ร้านนำร่อง 12 แห่งลดของเสียเฉลี่ยร้อยละ 45 ภายในแปดสัปดาห์ บทเรียนสำคัญคือควรทดสอบสมมติฐานด้านพฤติกรรมก่อนลงทุนพัฒนาเทคโนโลยีเต็มรูปแบบ
\end{examplebox}

\section{การให้เครดิตอย่างมีวุฒิภาวะ}
โครงการนวัตกรรมเป็นผลงานร่วม Mentee ควรระบุให้ชัดว่า \enquote{ทีมทำอะไร} และ \enquote{ตนรับผิดชอบอะไร} การให้เครดิตผู้อื่นไม่ได้ลดคุณค่าของ Mentee ตรงกันข้าม กลับแสดงความเข้าใจระบบและความสามารถในการทำงานร่วมกัน

\section{โครงการที่ไม่สำเร็จ}
ความล้มเหลวที่วิเคราะห์อย่างจริงจังมีคุณค่ามากกว่าความสำเร็จที่อธิบายไม่ได้ Mentor ควรสร้างพื้นที่ให้ Mentee แยกผลลัพธ์ออกจากตัวตน และถามว่า ข้อมูลใดทำให้ต้องเปลี่ยนทิศ ต้นทุนใดหลีกเลี่ยงได้ และระบบใดควรสร้างเพื่อไม่ให้ข้อผิดพลาดเกิดซ้ำ

\diagnostic{การตัดสินใจใดในโครงการนี้บอกวิธีคิดของ Mentee ได้ดีที่สุด และข้อมูลใดอยู่เบื้องหลังการตัดสินใจนั้น}

\begin{mentornote}
เมื่อเรื่องเล่าเต็มไปด้วยคำว่า \enquote{ได้ดำเนินการ} ให้ถามหา subject และ decision: ใครตัดสินใจอะไร เพราะเห็นข้อมูลใด
\end{mentornote}

\begin{keytakeaway}
Projects ที่ดีไม่ใช่รายการทุน แต่เป็นหลักฐานเชิงเรื่องเล่าที่ทำให้ผู้อ่านมองเห็นความสามารถ การตัดสินใจ ผลลัพธ์ และการเรียนรู้ของ Mentee
\end{keytakeaway}

\begin{reflection}
เลือกโครงการหนึ่งแล้วเขียนเรื่องเล่า 200 คำตามห้าจังหวะ ขีดเส้นใต้ทุกข้อกล่าวอ้างที่ยังไม่มี Evidence และวงกลมการตัดสินใจที่เป็นของ Mentee โดยตรง
\blankline\blankline
\end{reflection}

\chaptersummary{การคัดเลือกและเล่า Projects อย่างมีวินัยทำให้ประสบการณ์กลายเป็นหลักฐานของความสามารถ เรื่องเล่าที่ดีให้เครดิตทีม เปิดเผยการตัดสินใจ และไม่ซ่อนบทเรียนจากสิ่งที่ไม่เป็นไปตามแผน}
